\begin{definition}
For a real number $p$ put
\[
  d=d(p)=\frac16\Biggl(1+\frac1{p-1}+\sqrt{1+\frac{14}{p-1}+\frac1{(p-1)^2}}\,\Biggr).
\]
This is the exponent $d=d(p)$ of equation (50) of the paper, which Section~6
introduces as the smallest of the numbers $n/\gamma_n$, $n=1,2,\dots$

Two conventions make the formula total: $\sqrt t=0$ for $t<0$, and $1/0=0$.
The paper uses $d(p)$ only for $1<p<\infty$, where $p-1>0$ and the radicand is
positive, so neither convention is invoked.

The scan of the article prints the last denominator as $(p-2)^2$; the
transcription records that the correct denominator is $(p-1)^2$, and that is
the reading used here.  It is the reading for which $d=1/\gamma_1$ and for
which the right-hand side of equation (9) equals $1+d$.
\end{definition}
